\documentclass[10pt,a4paper]{amsart}
\usepackage[margin=19mm]{geometry}
\usepackage{amsmath,amssymb,mathtools}
\usepackage[scr=rsfs,cal=euler]{mathalfa}
\usepackage{xfrac}
\usepackage[unicode,hidelinks,bookmarks=false]{hyperref}
\newcommand{\ie}{i.e.\ }
\newcommand{\cf}{cf.\ }
\newcommand{\resp}{respectively\ }
\newcommand{\Subsection}{\subsection}
\providecommand{\nobreakdash}{\nobreak\hbox{-}}
\setlength{\emergencystretch}{2em}
\multlinegap0pt
\usepackage{bm}
\usepackage{amssymb}
\usepackage{mathtools}
\usepackage[shortlabels]{enumitem}
\usepackage{tikz}
\usepackage{algorithm}\usepackage{algpseudocode}
\usepackage{float}
\usepackage{xcolor}
\newtheorem{proposition}{Proposition}
\newtheorem{theorem}{Theorem}
\newcommand{\E}{\mathbb{E}}
\newcommand{\RKHS}{\mathcal{H}_k}
\newcommand{\Var}{\mathrm{Var}}
\newcommand{\argmin}{\operatornamewithlimits{arg\,min}}
\newcommand{\bI}{\mathbf{I}}
\newcommand{\bK}{\mathbf{K}}
\newcommand{\bS}{\mathbf{S}}
\title{\LARGE Kernel Methods for Stochastic Dynamical Systems with Application to Koopman Eigenfunctions:\\[0.5em] \large Feynman--Kac Representations and RKHS Approximation}
\author{Boumediene Hamzi, Houman Owhadi,, Umesh Vaidya}
\date{Source: arXiv:2603.01077v1 (CC BY 4.0)}
\begin{document}
\maketitle

\noindent\emph{Source and licence.} \LARGE Kernel Methods for Stochastic Dynamical Systems with Application to Koopman Eigenfunctions:\\[0.5em] \large Feynman--Kac Representations and RKHS Approximation. Version arXiv:2603.01077v1 (CC BY 4.0), distributed by arXiv under the Creative Commons Attribution 4.0 International licence, http://creativecommons.org/licenses/by/4.0/, as recorded in the arXiv licence metadata for this exact version and shown on the abstract page https://arxiv.org/abs/2603.01077v1. Full licence notice: \texttt{licenses/arxiv-2603.01077-CC-BY-4.0.txt}.\par\smallskip

\noindent\textit{Editorial scope.} Counted unit: the article's theorem thm:total\_error (source0.tex lines 1127--1132). The article's complete original proof of it is delivered with it (source0.tex lines 1134--1200, one \texttt{proof} environment of 67 lines). No mathematics is written, repaired or re-derived: the statement and the proof are the article's own lines, in the article's own notation, and no retained line is substituted or re-flowed.  Retained uncounted as the source-local context the proof needs: lines 880--886.  Nothing was omitted from any retained span, so the package carries no omission record.  No retained line contains a citation, so the wrapper carries no bibliography and no bibliography entry.  No retained line contains a figure environment, an image-inclusion command or drawing code, so no image is delivered and the package declares no asset dependency.  Editorial formatting only, in addition: an AMS \texttt{amsart} preamble that restates the article's own declarations and macros used by the retained lines, namely the environments \texttt{proposition}, \texttt{theorem} on separate counters, together with the standard packages those lines need.  The retained environments print in this document as Proposition 1, Theorem 1 in order of appearance, whereas the article prints the same items with the numbering of its own full document.  The complete native source of the article as distributed in the arXiv e-print for this exact version is bundled unchanged as \texttt{sources/195535/source0.tex}.
\begin{proposition}[Monte Carlo Error]\label{prop:mc_error}
The Monte Carlo estimator $\widehat{h}(x)$ satisfies
\begin{equation}
\E[(\widehat{h}(x) - h(x))^2] = \frac{\Var(h^{(1)}(x))}{K} \leq \frac{\sigma_{\mathrm{MC}}^2(x)}{K},
\end{equation}
where $\sigma_{\mathrm{MC}}^2(x)$ depends on the boundary data, source term, and exit time distribution.
\end{proposition}

\begin{theorem}[Total Error Bound]\label{thm:total_error}
Let $h^*$ be the true solution, $\{x_i\}_{i=1}^N$ collocation points, and $\widehat{h}(x_i)$ Monte Carlo estimates with $K$ samples. The kernel ridge regression estimator $\widehat{h}_N$ satisfies:
\begin{equation}
\frac{1}{N}\sum_{i=1}^N (\widehat{h}_N(x_i) - h^*(x_i))^2 \leq \underbrace{\inf_{u \in \RKHS}\left\{\frac{1}{N}\sum_{i=1}^N(u(x_i) - h^*(x_i))^2 + \eta\|u\|_{\RKHS}^2\right\}}_{\text{RKHS approximation error}} + \underbrace{\frac{C\sigma_{\mathrm{MC}}^2}{K}}_{\text{Monte Carlo error}}.
\end{equation}
\end{theorem}

\begin{proof}
We decompose the total error into approximation and estimation components.

\textbf{Step 1: Setup.}
Let $\mathbf{h}^* = (h^*(x_1), \ldots, h^*(x_N))^\top$ be the true values at collocation points, and $\widehat{\mathbf{h}} = (\widehat{h}(x_1), \ldots, \widehat{h}(x_N))^\top$ be the Monte Carlo estimates. The kernel ridge regression estimator solves:
\begin{equation}
\widehat{h}_N = \argmin_{u \in \RKHS} \frac{1}{N}\sum_{i=1}^N (u(x_i) - \widehat{h}(x_i))^2 + \eta \|u\|_{\RKHS}^2.
\end{equation}

\textbf{Step 2: Bias-variance decomposition.}
We write $\widehat{h}(x_i) = h^*(x_i) + \epsilon_i$, where $\epsilon_i = \widehat{h}(x_i) - h^*(x_i)$ is the Monte Carlo error at point $x_i$. By Proposition~\ref{prop:mc_error}:
\begin{equation}
\E[\epsilon_i] = 0, \quad \E[\epsilon_i^2] = \frac{\sigma_{\mathrm{MC}}^2(x_i)}{K} \leq \frac{\sigma_{\mathrm{MC}}^2}{K}.
\end{equation}

\textbf{Step 3: Oracle inequality.}
Define the ``oracle'' estimator $h_N^*$ as the kernel ridge regression fit to the true values:
\begin{equation}
h_N^* = \argmin_{u \in \RKHS} \frac{1}{N}\sum_{i=1}^N (u(x_i) - h^*(x_i))^2 + \eta \|u\|_{\RKHS}^2.
\end{equation}
The oracle satisfies the standard approximation bound:
\begin{equation}
\frac{1}{N}\sum_{i=1}^N (h_N^*(x_i) - h^*(x_i))^2 + \eta\|h_N^*\|_{\RKHS}^2 \leq \inf_{u \in \RKHS}\left\{\frac{1}{N}\sum_{i=1}^N(u(x_i) - h^*(x_i))^2 + \eta\|u\|_{\RKHS}^2\right\}.
\end{equation}

\textbf{Step 4: Perturbation analysis.}
The kernel ridge regression estimator has the closed form:
\begin{equation}
\widehat{h}_N(x) = \mathbf{k}(x)^\top (\bK + N\eta \bI)^{-1} \widehat{\mathbf{h}},
\end{equation}
where $\mathbf{k}(x) = (k(x, x_1), \ldots, k(x, x_N))^\top$. Similarly, the oracle is:
\begin{equation}
h_N^*(x) = \mathbf{k}(x)^\top (\bK + N\eta \bI)^{-1} \mathbf{h}^*.
\end{equation}
Therefore:
\begin{equation}
\widehat{h}_N(x) - h_N^*(x) = \mathbf{k}(x)^\top (\bK + N\eta \bI)^{-1} \boldsymbol{\epsilon},
\end{equation}
where $\boldsymbol{\epsilon} = \widehat{\mathbf{h}} - \mathbf{h}^* = (\epsilon_1, \ldots, \epsilon_N)^\top$.

\textbf{Step 5: Error decomposition.}
Using the triangle inequality:
\begin{align}
\frac{1}{N}\sum_{i=1}^N (\widehat{h}_N(x_i) - h^*(x_i))^2 &\leq \frac{2}{N}\sum_{i=1}^N (\widehat{h}_N(x_i) - h_N^*(x_i))^2 + \frac{2}{N}\sum_{i=1}^N (h_N^*(x_i) - h^*(x_i))^2.
\end{align}

\textbf{Step 6: Bounding the Monte Carlo contribution.}
Let $\bS = (\bK + N\eta \bI)^{-1} \bK$. Then:
\begin{align}
\frac{1}{N}\sum_{i=1}^N (\widehat{h}_N(x_i) - h_N^*(x_i))^2 &= \frac{1}{N}\boldsymbol{\epsilon}^\top \bS^\top \bS \boldsymbol{\epsilon} \leq \frac{\|\bS\|^2}{N} \|\boldsymbol{\epsilon}\|^2.
\end{align}
Taking expectations:
\begin{equation}
\E\left[\frac{\|\boldsymbol{\epsilon}\|^2}{N}\right] = \frac{1}{N}\sum_{i=1}^N \E[\epsilon_i^2] \leq \frac{\sigma_{\mathrm{MC}}^2}{K}.
\end{equation}
Since $\|\bS\| \leq 1$ (the smoother matrix has spectral norm at most 1), we obtain:
\begin{equation}
\E\left[\frac{1}{N}\sum_{i=1}^N (\widehat{h}_N(x_i) - h_N^*(x_i))^2\right] \leq \frac{\sigma_{\mathrm{MC}}^2}{K}.
\end{equation}

\textbf{Step 7: Combining bounds.}
Combining Steps 4--6:
\begin{align}
\E\left[\frac{1}{N}\sum_{i=1}^N (\widehat{h}_N(x_i) - h^*(x_i))^2\right] &\leq 2 \cdot \frac{\sigma_{\mathrm{MC}}^2}{K} + 2 \inf_{u \in \RKHS}\left\{\frac{1}{N}\sum_{i=1}^N(u(x_i) - h^*(x_i))^2 + \eta\|u\|_{\RKHS}^2\right\}.
\end{align}
Setting $C = 2$ yields the stated bound.
\end{proof}

\end{document}
